% ch8_b §8.3–8.4 (书页 266–271 / p280–p285)
%--A-TAIL--
%JOIN%
介质，而是占据格子中的一个格点；在那里，它被“其他”同样极化了的原子所包围，却并不被它们贯穿（见 §2.3）。众所周知，在初等静电学和静磁学中，局域场必须有一个修正；如果包围原子的“空腔”多少近似球形，那么我们就得到局域场的 \emph{Lorentz 修正}（Lorentz correction）
\begin{equation*}
\mathbf{E}_{\mathrm{local}} = \mathbf{E}_{\mathrm{macroscopic}} + \frac{4}{3}\pi\mathbf{P},\tag{8.52}
\end{equation*}
其中 $\mathbf{P}$ 是平均极化。这就导致 \emph{Clausius–Mosotti 关系}（Clausius–Mosotti relation）
\begin{equation*}
\left(\frac{\epsilon-1}{\epsilon+2}\right) = \frac{4}{3}\pi N\alpha\tag{8.53}
\end{equation*}
以取代 (8.42)。结果是，观测到的介电常数和折射率将遵循一个比 (8.51) 更复杂的公式；例如，“完全”反射的频段就可能得以避免。

当介电极化是由于液体或固体分子上\emph{永久}偶极矩（permanent dipole moments）的排齐所致时，我们必须引入进一步的修正，例如 \emph{Onsager 局域场}（Onsager local field）模型，它把邻近原子矩取向的涨落考虑了进去。

\section{离子晶体中的光学模}

让我们回到格子的运动方程，如 §2.1 中那样。那里已经证明，第 $l$ 个晶胞中格点 $s$ 上的离子，其位移满足 (2.7)，即
\begin{equation*}
M_s \ddot{\mathbf{u}}_{sl} = -\sum_{s'\boldsymbol{h}} \mathsf{G}_{ss'}(\boldsymbol{h})\cdot\mathbf{u}_{s',\,l+\boldsymbol{h}}.\tag{8.54}
\end{equation*}
现在假设我们用形如 (8.3) 的电磁波来驱动这个系统。必须在 (8.54) 的右端加上一个力
\begin{equation*}
\boldsymbol{e}_s\,\mathbf{E}_0\exp\{i(\mathbf{K}\cdot\boldsymbol{l}-\omega t)\},\tag{8.55}
\end{equation*}
其中 $\boldsymbol{e}_s$ 是晶胞中格点 $s$ 上的离子所带的电荷。

要处理这样的非齐次项，我们只需要微分方程 (8.54) 的一个特解。显然，采用 (2.8) 的假设并令 $\boldsymbol{q}=\mathbf{K}$，随空间变化的因子即被消去；我们便得到波矢为 $\mathbf{K}$ 的模的振幅 $\mathbf{U}_{s\mathbf{K}}$ 所满足的方程
\begin{equation*}
M_s \ddot{\mathbf{U}}_{s\mathbf{K}} = -\sum \mathsf{G}_{ss'}(\mathbf{K})\cdot\mathbf{U}_{s'\mathbf{K}} + \boldsymbol{e}_s\,\mathbf{E}_0\,e^{-i\omega t}.\tag{8.56}
\end{equation*}
要看出这些方程之解的形式，让我们考虑 (2.22) 中那个双原子线性链的简单模型。实际上，我们所考虑的电磁波在原子的尺度上波长非常之大，
因此可以取 $\cos Ka \approx 1$。如果两个原子带有异号电荷——在典型的离子晶体中情形正是如此——那么我们就有
\begin{equation*}
\left.
\begin{aligned}
M_1 \ddot{U}_1 &= 2\alpha(U_2-U_1) + \boldsymbol{e}E_0\,e^{-i\omega t},\\
M_2 \ddot{U}_2 &= 2\alpha(U_1-U_2) - \boldsymbol{e}E_0\,e^{-i\omega t}.
\end{aligned}
\right\}
\tag{8.57}
\end{equation*}
这两个方程有解
\begin{equation*}
(U_1-U_2) = \frac{\boldsymbol{e}E_0(1/M_1+1/M_2)}{\omega_T^2-\omega^2}\,e^{-i\omega t},\tag{8.58}
\end{equation*}
其中 $\omega_T$ 是系统在 $\boldsymbol{q}=0$ 处的固有频率，即由 (2.25) 给出的值。我们可以算出与这种运动相联系的偶极矩，并把它表示成晶格的一个极化率
\begin{equation*}
4\pi N\alpha(\omega) = \frac{4\pi N\boldsymbol{e}^2(1/M_1+1/M_2)}{\omega_T^2-\omega^2}.\tag{8.59}
\end{equation*}
于是，一个纯粹经典的论证就把我们引到了一个类似于 (8.42) 的色散公式。(8.59) 的分子可以写成 $\Omega_p^{\prime 2}$，表明它非常像离子集合体的等离体频率 (6.81)。由于离子比电子重得多，它必定远小于 $\omega_p^2$。另一方面，$\omega_T$ 正是长波长晶格振动的频率，比原子或离子中任何电子跃迁的能量都小得多。因此，在很宽的频率范围内，我们可以写成
\begin{equation*}
\epsilon(\omega) \approx \epsilon(\infty) + \frac{\Omega_p^{\prime 2}}{\omega_T^2-\omega^2},\tag{8.60}
\end{equation*}
其中误差来自 (8.53) 那样的局域场效应。在这个公式中，$\epsilon(\infty)$ 表示在 $\omega \gg \omega_T$（即电子跃迁尚未被激发之前）由折射率的观测推出的介电常数。由此我们可以导出静态介电常数的 \emph{Born 方程}（Born equation）
\begin{equation*}
\epsilon(0) \approx \epsilon(\infty) + \Omega_p^{\prime 2}/\omega_T^2.\tag{8.61}
\end{equation*}
在三维离子晶体中，$\omega_T$ 将是一个\emph{横向}极化的光学支的频率，其偶极矩与横向电磁波强烈地相互作用。实际上，我们这里得到的是一个\emph{极化激元}（polariton）——光与晶体极化的一种混合模，它才是材料中真正的传播模。当 $\omega$ 扫过 $\omega_T$ 时，介电常数变为负值：正如 (5.67) 和 (8.46) 的情形，辐射将被晶体表面完全反射。\emph{剩余射线效应}（Reststrahlen effect）——晶体中电磁传播的一个禁区——将一直持续到譬如说 $\omega_L$ 处，在该处 (8.60) 重新变为正，即
\begin{equation*}
\omega_L^2 = \omega_T^2 + \Omega_p^{\prime 2}/\epsilon(\infty).\tag{8.62}
\end{equation*}
但介电常数的定义意味着（参见 (5.59)）：在这个频率上，无需任何外部激发，晶体内部就可以存在非常大的电场变化。这正是离子晶格中等离体振荡的类比，其中既有局域原子间力的贡献，也有长程静电场的贡献：$\omega_L$ 就是晶格谱中\emph{纵向}光学支的长波极限。由 (8.61) 和 (8.62) 我们得到与 (8.60) 相配的一个关系——\emph{Lyddane–Sachs–Teller 关系}（Lyddane–Sachs–Teller relation）
\begin{equation*}
\omega_L^2/\omega_T^2 = \epsilon(0)/\epsilon(\infty)\tag{8.63}
\end{equation*}
它对任何极性材料都相当普遍地成立。

实际上，(8.59) 和 (8.60) 的分母中应当含有一个虚项（参见 (8.51)），对应于耗散效应。例如，通过产生声子，电磁场的能量可能被吸收，而这些声子将通过力常数的非谐性（anharmonicity）与光学支相联系（§§2.11, 7.10, 8.4）。这类效应将是依赖于温度的，并且也应能被观察到，表现为谐和近似 (2.25) 或 (8.53) 中有效力常数 $\alpha$ 的变化。甚至还可能出现这样的情况：对某一组晶格位移而言，原子间力的平衡如此精妙，以致这种温度依赖使 $\omega_T^2$ 在某个临界温度 $T_c$ 以下取负值——这个事件我们在宏观上应当认作是向\emph{铁电相}（ferro-electric phase）的转变（§10.5），晶格结构从此带有永久的、内建的介电极化。

纵向光学支的频率比较低，这一点在\emph{极化子}（polaron）理论中很重要。离子晶体中的一个电子携带着静电场，使其周围的晶格极化。对于静止的电子，我们本可以用介电常数 $\epsilon(0)$；但晶格对运动电子的响应在动力学上受频率 $\omega_L$ 的限制，它描写介质介电极化的固有振荡。

这种情形与等离激元理论（§5.8）中的情形非常相似。如果电子的能量超过 $\hbar\omega_L$，就会产生一个光学声子，电子将被减速并受到强烈散射。即使电子运动得很慢，我们也必须考虑光学模中量子的虚激发与再吸收；
换言之，电子的行为就像一个被一团虚声子云包围着的粒子。这个复合客体的自能是速度的函数，因而载流子的有效“质量”会显著增大。

\emph{大极化子}（large polaron）的理想化哈密顿量，作为一个强相互作用的场论系统的模型，已被广泛研究过，但它在实践中并不是一个容易观察的对象。这个模型本身只在极化区域比晶格的一个晶胞多少大一些的时候才有效；否则，更合适的是自陷的\emph{小极化子}（small polaron）模型（§6.7）。

\section{光子–声子跃迁}

用场论的语言来说，晶体对红外的吸收被描写为一个\emph{光子}（photon）与一个或多个声子的相互作用。这类过程的选择定则，已经在 §§2.7、2.8 中以\emph{非弹性衍射}（inelastic diffraction）的名义导出过。例如在最简单的情形，衍射束波矢的\emph{改变}必须等于所发射声子的波矢，即
\begin{equation*}
\mathbf{k}-\mathbf{k}' = \mathbf{q}\tag{8.64}
\end{equation*}
同时辐射频率的改变必须等于声子的频率，
\begin{equation*}
\omega-\omega' = \nu\tag{8.65}
\end{equation*}
如 (2.101) 和 (2.102) 中那样。在某些特殊情形下，可以满足这些条件而不产生任何出射的衍射波；光子的全部能量和动量都转移给了晶体。

这里要注意的要点是，我们处理的是可见和红外波段的光，其波长比晶体的晶格常数大得多。因此，$|\mathbf{k}-\mathbf{k}'|$ 在布里渊区的尺度上必定极其微小。对于一切寻常的光学现象，我们可以取 $\mathbf{k}-\mathbf{k}'\approx 0$，而把注意力集中在区图式所允许的能量改变上；这一方面恰好与适用于 X 射线和中子的“衍射”观点互为补充——在后者中，动量的改变很大，但能量差却不容易探测。

我们这里考虑的过程，即一个光学声子的吸收或发射，因而在简约区图式中由 $\mathbf{q}=0$ 处一条近乎竖直的线来表示。如果光学频谱有好几个分支，那么原则上应当能观察到好几条不同的谱线，如图 141(b) 那样。但是这类过程的
强度将取决于电磁波与晶体耦合的矩阵元。正如 §8.3 中所见，对于大多数极性晶体的横向光学支，这种耦合强到如此程度，以致用弱相互作用的光子和声子来描述已不恰当，我们得到的将是剩余射线效应。但在某些化合物半导体中，可以找到偶极矩非常弱的光学支，在透过薄膜的透射光中能够观察到\emph{单声子吸收}（one-phonon absorption）。

%FIG{141}: {光学模的吸收：(a) 一维图式；(b) 三维图式。}

观察到光子的吸收（和发射）过程在布里渊区中是“竖直”的，这提示我们还可能有别的“竖直”跃迁，如图 142 所示。这是一个涉及一个声学支和一个光学支的跃迁。被吸收的光子，其能量显然由下式构成
\begin{equation*}
\hbar\omega = \hbar\nu_{\mathrm{optical}} - \hbar\nu_{\mathrm{acoustic}}.\tag{8.66}
\end{equation*}
实际上就是：一个光子和一个声学声子消失了；一个光学声子产生了。动量条件也同样得到满足。我们可以忽略光子的动量，简单地断言
\begin{equation*}
\hbar\mathbf{q}_{\mathrm{optical}} = \hbar\mathbf{q}_{\mathrm{acoustic}},\tag{8.67}
\end{equation*}
这正是“跃迁必须竖直”这句话所蕴含的意思。

%FIG{142}: {多声子吸收。}

这类过程为选择定则所允许，因此可以观察到。它们在远红外区给出整段整段的吸收带。这种吸收带的宽度和总体结构，显然
取决于从声学支到光学支可能发生的跃迁的种类及其间距。对这类过程作完整的讨论显然相当复杂，因为我们处理的是两个函数之差，而其中每一个函数都是布里渊区中 $\mathbf{q}$ 的函数。

这还没有穷尽一切可能的过程。我们可以把上述跃迁称为“差式”（difference）过程。同样还可能有\emph{求和带}（summation bands），其中由光子的能量产生\emph{两个}声子，但它们的动量大小相等、方向相反。在那种情形下
\begin{equation*}
\hbar\omega = \hbar\nu_{\mathrm{optical}} + \hbar\nu_{\mathrm{acoustic}},\tag{8.68}
\end{equation*}
并有
\begin{equation*}
\hbar\mathbf{q}_{\mathrm{optical}} = -\hbar\mathbf{q}_{\mathrm{acoustic}}.
\end{equation*}
要真正计算跃迁概率，我们需要知道偶极矩作为晶格位移的函数中的非线性项。例如，我们需要在如下表达式中的张量 $\mathbf{B}_{ss'}$
\begin{equation*}
\mathbf{M}_l = \sum_s \mathbf{A}_s\cdot\mathbf{u}_{sl} + \sum_{ss'} \mathbf{u}_{sl}\cdot\mathbf{B}_{ss'}\cdot\mathbf{u}_{s'l}\tag{8.69}
\end{equation*}
它给出第 $l$ 个晶胞中由位移 $\mathbf{u}_{sl}$ 产生的偶极矩。于是我们用 (2.8)、(2.10) 和 (2.129) 把每一位移用谱的各个分支的声子的湮没和产生算符表出。$\mathbf{B}_{ss'}$ 项对应于两个此类算符的乘积，它给出我们所寻求的双声子过程。跃迁的实际概率将包含如下形式的两个矩阵元的平方
\begin{equation*}
\left.
\begin{gathered}
\langle n_{\mathbf{q}}+1|a_{\mathbf{q}}^{*}|n_{\mathbf{q}}\rangle = (n_{\mathbf{q}}+1)^{\frac{1}{2}},\\
\text{或}\quad \langle n_{\mathbf{q}}-1|a_{\mathbf{q}}|n_{\mathbf{q}}\rangle = n_{\mathbf{q}}^{\frac{1}{2}},
\end{gathered}
\right\}
\tag{8.70}
\end{equation*}
对应于声子态的占据数增加或减少一个量子。由于 $n_{\mathbf{q}}$ 的平均值是温度的函数（如 (2.46) 那样），这些过程将是依赖于温度的。不过一般说来，非线性系数 $\mathbf{B}_{ss'}$ 很小，所以与单声子过程相比，这些带是弱的。

值得指出的是，对于任何一个分支，函数 $\nu_{\mathbf{q}}$ 都是 $\mathbf{q}$ 的连续周期函数，其周期为倒格子的周期（参见 §2.2）。由此可知，任何一对这类函数的和或差也具有同样的性质。于是，van Hove 定理（§2.5）适用于频率的分布；在观测谱中只可能出现图 27 所示的那四种类型的奇点（除非矩阵元引入额外的奇点）。
